% Options for packages loaded elsewhere
% Options for packages loaded elsewhere
\PassOptionsToPackage{unicode}{hyperref}
\PassOptionsToPackage{hyphens}{url}
%
\documentclass[
  english,
  russian,
  12pt,
  a4paper,
  DIV=11,
  numbers=noendperiod]{scrreprt}
\usepackage{xcolor}
\usepackage{amsmath,amssymb}
\setcounter{secnumdepth}{5}
\usepackage{iftex}
\ifPDFTeX
  \usepackage[T1]{fontenc}
  \usepackage[utf8]{inputenc}
  \usepackage{textcomp} % provide euro and other symbols
\else % if luatex or xetex
  \usepackage{unicode-math} % this also loads fontspec
  \defaultfontfeatures{Scale=MatchLowercase}
  \defaultfontfeatures[\rmfamily]{Ligatures=TeX,Scale=1}
\fi
\usepackage{lmodern}
\ifPDFTeX\else
  % xetex/luatex font selection
\fi
% Use upquote if available, for straight quotes in verbatim environments
\IfFileExists{upquote.sty}{\usepackage{upquote}}{}
\IfFileExists{microtype.sty}{% use microtype if available
  \usepackage[]{microtype}
  \UseMicrotypeSet[protrusion]{basicmath} % disable protrusion for tt fonts
}{}
\usepackage{setspace}
% Make \paragraph and \subparagraph free-standing
\makeatletter
\ifx\paragraph\undefined\else
  \let\oldparagraph\paragraph
  \renewcommand{\paragraph}{
    \@ifstar
      \xxxParagraphStar
      \xxxParagraphNoStar
  }
  \newcommand{\xxxParagraphStar}[1]{\oldparagraph*{#1}\mbox{}}
  \newcommand{\xxxParagraphNoStar}[1]{\oldparagraph{#1}\mbox{}}
\fi
\ifx\subparagraph\undefined\else
  \let\oldsubparagraph\subparagraph
  \renewcommand{\subparagraph}{
    \@ifstar
      \xxxSubParagraphStar
      \xxxSubParagraphNoStar
  }
  \newcommand{\xxxSubParagraphStar}[1]{\oldsubparagraph*{#1}\mbox{}}
  \newcommand{\xxxSubParagraphNoStar}[1]{\oldsubparagraph{#1}\mbox{}}
\fi
\makeatother


\usepackage{longtable,booktabs,array}
\usepackage{calc} % for calculating minipage widths
% Correct order of tables after \paragraph or \subparagraph
\usepackage{etoolbox}
\makeatletter
\patchcmd\longtable{\par}{\if@noskipsec\mbox{}\fi\par}{}{}
\makeatother
% Allow footnotes in longtable head/foot
\IfFileExists{footnotehyper.sty}{\usepackage{footnotehyper}}{\usepackage{footnote}}
\makesavenoteenv{longtable}
\usepackage{graphicx}
\makeatletter
\newsavebox\pandoc@box
\newcommand*\pandocbounded[1]{% scales image to fit in text height/width
  \sbox\pandoc@box{#1}%
  \Gscale@div\@tempa{\textheight}{\dimexpr\ht\pandoc@box+\dp\pandoc@box\relax}%
  \Gscale@div\@tempb{\linewidth}{\wd\pandoc@box}%
  \ifdim\@tempb\p@<\@tempa\p@\let\@tempa\@tempb\fi% select the smaller of both
  \ifdim\@tempa\p@<\p@\scalebox{\@tempa}{\usebox\pandoc@box}%
  \else\usebox{\pandoc@box}%
  \fi%
}
% Set default figure placement to htbp
\def\fps@figure{htbp}
\makeatother



\ifLuaTeX
\usepackage[bidi=basic,provide=*]{babel}
\else
\usepackage[bidi=default,provide=*]{babel}
\fi
% get rid of language-specific shorthands (see #6817):
\let\LanguageShortHands\languageshorthands
\def\languageshorthands#1{}


\setlength{\emergencystretch}{3em} % prevent overfull lines

\providecommand{\tightlist}{%
  \setlength{\itemsep}{0pt}\setlength{\parskip}{0pt}}



 
\usepackage[backend=biber,langhook=extras,autolang=other*]{biblatex}
\addbibresource{bib/cite.bib}

\usepackage[]{csquotes}

\usepackage{indentfirst}
\usepackage{float}
\floatplacement{figure}{H}
\IfFileExists{plex-otf.sty}{
  %% Full TeXlive
  \IfPackageAtLeastTF{plex-otf.sty}{2024-12-06}{
    \usepackage[math,RM={Scale=0.94},SS={Scale=0.94},SScon={Scale=0.94},TT={Scale=MatchLowercase,FakeStretch=0.9},DefaultFeatures={Ligatures=Common}]{plex-otf}
  }{
    \usepackage[RM={Scale=0.94},SS={Scale=0.94},SScon={Scale=0.94},TT={Scale=MatchLowercase,FakeStretch=0.9},DefaultFeatures={Ligatures=Common}]{plex-otf}
  }
}{
  %% TinyTeX
  \usepackage{libertine}
}
\KOMAoption{captions}{tableheading}
\makeatletter
\@ifpackageloaded{caption}{}{\usepackage{caption}}
\AtBeginDocument{%
\ifdefined\contentsname
  \renewcommand*\contentsname{Содержание}
\else
  \newcommand\contentsname{Содержание}
\fi
\ifdefined\listfigurename
  \renewcommand*\listfigurename{Список иллюстраций}
\else
  \newcommand\listfigurename{Список иллюстраций}
\fi
\ifdefined\listtablename
  \renewcommand*\listtablename{Список таблиц}
\else
  \newcommand\listtablename{Список таблиц}
\fi
\ifdefined\figurename
  \renewcommand*\figurename{Рисунок}
\else
  \newcommand\figurename{Рисунок}
\fi
\ifdefined\tablename
  \renewcommand*\tablename{Таблица}
\else
  \newcommand\tablename{Таблица}
\fi
}
\@ifpackageloaded{float}{}{\usepackage{float}}
\floatstyle{ruled}
\@ifundefined{c@chapter}{\newfloat{codelisting}{h}{lop}}{\newfloat{codelisting}{h}{lop}[chapter]}
\floatname{codelisting}{Список}
\newcommand*\listoflistings{\listof{codelisting}{Листинги}}
\makeatother
\makeatletter
\makeatother
\makeatletter
\@ifpackageloaded{caption}{}{\usepackage{caption}}
\@ifpackageloaded{subcaption}{}{\usepackage{subcaption}}
\makeatother
\usepackage{bookmark}
\IfFileExists{xurl.sty}{\usepackage{xurl}}{} % add URL line breaks if available
\urlstyle{same}
\hypersetup{
  pdftitle={Шаблон отчёта по лабораторной работе},
  pdfauthor={Дмитрий Сергеевич Кулябов},
  pdflang={ru-RU},
  hidelinks,
  pdfcreator={LaTeX via pandoc}}


\title{Шаблон отчёта по лабораторной работе}
\usepackage{etoolbox}
\makeatletter
\providecommand{\subtitle}[1]{% add subtitle to \maketitle
  \apptocmd{\@title}{\par {\large #1 \par}}{}{}
}
\makeatother
\subtitle{Простейший вариант}
\author{Дмитрий Сергеевич Кулябов}
\date{}
\begin{document}
\maketitle

\renewcommand*\contentsname{Содержание}
{
\setcounter{tocdepth}{1}
\tableofcontents
}
\listoffigures
\listoftables

\setstretch{1.5}
\chapter{Цель
работы}\label{ux446ux435ux43bux44c-ux440ux430ux431ux43eux442ux44b}

Приобретение практических навыков по конфигурированию SMTP-сервера в
части настройки аутентификации

\chapter{Выполнение лабораторной
работы}\label{ux432ux44bux43fux43eux43bux43dux435ux43dux438ux435-ux43bux430ux431ux43eux440ux430ux442ux43eux440ux43dux43eux439-ux440ux430ux431ux43eux442ux44b}

Во втором терминале сервера запустим вывод логов почты (рис.
\autocite*{fig:001}).

\begin{figure}

{\centering \pandocbounded{\includegraphics[keepaspectratio]{image/1.png}}

}

\caption{Логи почты /var/log/maillog}

\end{figure}%

В первом же терминале сервера авторизуемся под рутом и откроем файл
конфигурации /etc/dovecot/dovecot.conf.

В этом файле допишем в список рабочих протоколов протокол lmtp (рис.
\autocite*{fig:003}).

\begin{figure}

{\centering \pandocbounded{\includegraphics[keepaspectratio]{image/3.png}}

}

\caption{/etc/dovecot/dovecot.conf}

\end{figure}%

Далее, откроем конфигурационный файл /etc/dovecot/conf.d/10-master.conf.
Внутри этого файла пропишем следующее тело для структуры service lmtp
(рис. \autocite*{fig:005}).

\begin{figure}

{\centering \pandocbounded{\includegraphics[keepaspectratio]{image/5.png}}

}

\caption{service lmtp}

\end{figure}%

Далее, пропишем в postfix сокет, через который будет идти отправка
сообщений. После этого откроем файл /etc/dovecot/conf.d/10-auth.conf
(рис. \autocite*{fig:006}).

\begin{figure}

{\centering \pandocbounded{\includegraphics[keepaspectratio]{image/6.png}}

}

\caption{Настройка сокета}

\end{figure}%

В этом файле зададим значение для поля auth\_username\_format,
отвечающее за формат имени пользователя для аутентификации. В нашем
случае домен не будет указываться (рис. \autocite*{fig:007}).

\begin{figure}

{\centering \pandocbounded{\includegraphics[keepaspectratio]{image/7.png}}

}

\caption{auth\_username\_format}

\end{figure}%

Перезапустим postfix и dovecot (рис. \autocite*{fig:008}).

\begin{figure}

{\centering \pandocbounded{\includegraphics[keepaspectratio]{image/8.png}}

}

\caption{Перезапуск postfix и dovecot}

\end{figure}%

Теперь перейдём на виртуальную машину клиента. Попробуем отправить
письмо самому себе (рис. \autocite*{fig:009}).

\begin{figure}

{\centering \pandocbounded{\includegraphics[keepaspectratio]{image/9.png}}

}

\caption{Отправка письма}

\end{figure}%

В логах, которые мы открывали на сервере в самом начале выполнения
лабораторной работы, мы видим, что письмо было доставлено в ящик. Об
этом свидетельствует подпись \enquote{saved mail to INBOX}. Кроме того,
теперь в логах пишется, что транспортировка осуществляется через lmtp
(passing \ldots{} to transport=lmtp) (рис. \autocite*{fig:010}).

\begin{figure}

{\centering \pandocbounded{\includegraphics[keepaspectratio]{image/1.png}}

}

\caption{Логи почты}

\end{figure}%

Откроем на сервере почтовый ящик, чтобы убедится, что письмо успешно
доставлено. Как видим, это действительно так (рис. \autocite*{fig:011}).

\begin{figure}

{\centering \pandocbounded{\includegraphics[keepaspectratio]{image/11.png}}

}

\caption{Почтовый ящик mail}

\end{figure}%

Теперь откроем файл конфигурации по пути
/etc/dovecot/conf.d/10-master.conf.

Приведем содержание тела структуры service auth к следующему виду.
Разберём построчно

\begin{verbatim}
service auth { ... } - Эта строка объявляет начало секции конфигурации для внутренней службы Dovecot, которая называется auth.   

unix_listener /var/spool/postfix/private/auth { ... } - Указывает Dovecot создать "слушателя" на основе UNIX-сокета.   

group = postfix - Устанавливает группу-владельца для файла сокета.   
user = postfix - Устанавливает пользователя-владельца для файла сокета.   
mode = 0660 - Устанавливает права доступа к файлу сокета в восьмеричном формате.   
      
unix_listener auth-userdb { ... } - Создает второй UNIX-сокет.   

mode = 0600 - Устанавливает права доступа для этого внутреннего сокета.   

user = dovecot - Устанавливает пользователя-владельца dovecot.   
\end{verbatim}

(рис. \autocite*{fig:013}).

\begin{figure}

{\centering \pandocbounded{\includegraphics[keepaspectratio]{image/13.png}}

}

\caption{Структура service auth}

\end{figure}%

Теперь настроим аутентификацию почты smtp для postfix и укажем, каким
правилам следовать для работы с почтой и её фильтрации. Рассмотрим опции

reject\_unknown\_recipient\_domain - Отклонить письмо, если домен в
адресе получателя не существует.

permit\_mynetworks - Разрешить письмо без дальнейших проверок, если
IP-адрес клиента, отправляющего письмо, находится в списке доверенных
сетей.

reject\_non\_fqdn\_recipient - Отклонить письмо, если адрес получателя
не является полностью определённым доменным именем.

reject\_unauth\_destination - отклоняет письмо, если домен получателя не
является локальным для этого сервера и при этом сессия не
аутентифицирована.

reject\_unverified\_recipient - Отклонить письмо, если Postfix не может
проверить существование получателя.

permit - Если ни одно из предыдущих правил не отклонило и не разрешило
письмо, это правило разрешает его.

(рис. \autocite*{fig:014}).

\begin{figure}

{\centering \pandocbounded{\includegraphics[keepaspectratio]{image/14.png}}

}

\caption{Настройка почты в postfix}

\end{figure}%

Теперь отредактируем файл /etc/postfix/master.cf.

В этом файле внесём изменения так, чтобы smtp поддерживал авторизацию по
sasl (рис. \autocite*{fig:016}).

\begin{figure}

{\centering \pandocbounded{\includegraphics[keepaspectratio]{image/16.png}}

}

\caption{Включение авторизации}

\end{figure}%

Теперь перезапустим postfix и dovecot (рис. \autocite*{fig:017}).

\begin{figure}

{\centering \pandocbounded{\includegraphics[keepaspectratio]{image/8.png}}

}

\caption{Перезапуск postfix и dovecot}

\end{figure}%

Теперь на клиенте установим telnet (рис. \autocite*{fig:018}).

\begin{figure}

{\centering \pandocbounded{\includegraphics[keepaspectratio]{image/18.png}}

}

\caption{Установка Telnet}

\end{figure}%

Теперь получим ключ авторизации. Этот ключ представляет из себя строку,
содержащую имя пользователя и пароль, и зашифрованную в base64. Теперь
по telnet подключимся к почтовому серверу и проверим соединение. После
этого попробуем с помощью команды auth авторизироваться, в качестве
ключа используя нашу base64 строку. Как видим, авторизация прошла
успешно (Authentification successful) (рис. \autocite*{fig:019}).

\begin{figure}

{\centering \pandocbounded{\includegraphics[keepaspectratio]{image/19.png}}

}

\caption{Авторизация по telnet}

\end{figure}%

Теперь настроим сертификаты для postfix, а также уровень security и путь
к базе данных кэша (рис. \autocite*{fig:020}).

\begin{figure}

{\centering \pandocbounded{\includegraphics[keepaspectratio]{image/20.png}}

}

\caption{Настройка postfix}

\end{figure}%

Вновь откроем файл /etc/postfix/master.cf и изменим его следующим
образом (рис. \autocite*{fig:021}).

\begin{figure}

{\centering \pandocbounded{\includegraphics[keepaspectratio]{image/21.png}}

}

\caption{/etc/postfix/master.cf}

\end{figure}%

Теперь настроим firewall, разрешив использовать smtp-submission, и
перезапустим postfix (рис. \autocite*{fig:022}).

\begin{figure}

{\centering \pandocbounded{\includegraphics[keepaspectratio]{image/22.png}}

}

\caption{Настройка firewall}

\end{figure}%

Теперь подключимся к серверу через openssl (рис. \autocite*{fig:023}).

\begin{figure}

{\centering \pandocbounded{\includegraphics[keepaspectratio]{image/23.png}}

}

\caption{openssl}

\end{figure}%

Теперь сохраним внесённые нами изменения в vagrant (рис.
\autocite*{fig:029}).

\begin{figure}

{\centering \pandocbounded{\includegraphics[keepaspectratio]{image/29.png}}

}

\caption{vagrant}

\end{figure}%

На сервере изменим скрипт mail.sh следующим образом (рис.
\autocite*{fig:030}).

\begin{figure}

{\centering \pandocbounded{\includegraphics[keepaspectratio]{image/30.png}}

}

\caption{mail.sh для сервера}

\end{figure}%

\chapter{Выводы}\label{ux432ux44bux432ux43eux434ux44b}

В результате выполнения лабораторной работы были получены навыки
продвинутой настройки smtp и авторизации

\printbibliography[heading=none]





\end{document}
